\section{Reproducibility}\label{sec:repro}

\paragraph{Source repository.}
\url{https://github.com/roastercode/beamfs}, sub-directory
\texttt{papers/2026-04-beamfs-v1/}.

\paragraph{Companion report.}
The FTRFS v1 Technical Report~\cite{desbrieres2026ftrfs} is
prerequisite reading. Its empirical dataset is reused in
\cref{sec:empirical}; its threat model is extended in
\cref{sec:phys}.

\paragraph{Build commands.}
\begin{lstlisting}[language=bash]
cd papers/2026-04-beamfs-v1
make figures   # regenerate matplotlib figures from canonical data
make           # build paper.pdf via lualatex + biber + lualatex
\end{lstlisting}

\paragraph{Figure data generation.}
The three figures (\cref{fig:see-cross-section,fig:nand-retention,fig:recovery-surface}) are generated from canonical
cross-section curves and JEDEC NAND endurance data by Python scripts
in \texttt{scripts/}:
\begin{itemize}
\item \texttt{scripts/gen\_data.py} produces \texttt{data/fig1.dat}
  (Weibull fits to Petersen-form SEE cross-sections),
  \texttt{data/fig2.dat} (TLC NAND raw BER versus retention age,
  empirical Cai/Mielke model), and \texttt{data/fig3.dat}
  (RS-capacity-bounded recovery probability over the
  LET\(\times\)fluence grid, Poisson statistics).
\item \texttt{scripts/gen\_figures.py} reads these data and produces
  \texttt{figures/fig1.pdf}, \texttt{figures/fig2.pdf},
  \texttt{figures/fig3.pdf} via matplotlib.
\end{itemize}
Each script is deterministic: given the canonical input parameters
documented in its header, it produces a byte-for-byte identical
output.

\paragraph{Verification.}
\begin{lstlisting}[language=bash]
md5sum data/*.dat
# values are committed in scripts/gen_data.py headers
# matching values mean: same parameters, same output.
make
# zero error, zero warning expected.
\end{lstlisting}
